\section*{2. Grupoides cofibrados prorepresentables}

Sea $\underline A$ un grupoide cofibrado sobre
$\underline C_\Lambda$ que admite un objeto formal mínimamente versal
$\xi$. En general, $\xi$ por sí solo no determina $\underline A$
salvo $\underline C_\Lambda$-equivalencia. También es necesario saber
cuándo dos morfismos
\[
 R\mathrel{\substack{\xrightarrow{\ f\ }\\[-.6ex]
                      \xrightarrow[\ g\ ]{}}}S
\]
dan lugar a un $\underline C_\Lambda$-isomorfismo
$\xi(f)\longrightarrow\xi(g)$, donde $R=P(\xi)$. Como aplicación de
1.11 y 1.12, consideremos ahora el problema de recuperar $\underline A$
salvo $\underline C_\Lambda$-equivalencia. Obtenemos una respuesta
satisfactoria para una clase amplia de grupoides cofibrados.

%%%%%%%%%%  scan idx 61  |  folio 50  %%%%%%%%%%

Muchos grupoides cofibrados que aparecen en la práctica satisfacen una
propiedad ligeramente más fuerte que la semihomogeneidad; se definirá
explícitamente en 2.5. Para estudiarlos, conviene reformular la
discusión anterior en términos de «funtores lisos». Durante toda esta
sección siguen vigentes la notación y la terminología de la sección
anterior.

\underline{Definición 2.1.} Sea
$\varphi:\underline A\longrightarrow\underline B$ un funtor de
grupoides cofibrados sobre $\underline C_\Lambda$. Decimos que
$\varphi$ es \underline{liso} si todo morfismo sobreyectivo
$\beta:b'\longrightarrow\varphi(a)$ en $\underline B$ puede
completarse en un diagrama conmutativo
\begin{equation}\tag{*}
\begin{tikzcd}[row sep=large,column sep=huge]
 \varphi(a') \arrow[r] \arrow[dr,"{\varphi(\alpha)}"']
 & b' \arrow[d,"\beta"]\\
 & \varphi(a)
\end{tikzcd}
\end{equation}
para algún $\alpha:a'\longrightarrow a$ en $\underline A$. Decimos
que $\varphi$ es \underline{mínimamente liso} si es liso y
\[
 \varphi(k[\varepsilon]):
 [\underline A(k[\varepsilon])]
 \longrightarrow
 [\underline B(k[\varepsilon])]
\]
es biyectivo.

\underline{Observación 2.2(a).} Sea
$\varphi:\underline A\longrightarrow\underline B$ un funtor de
grupoides cofibrados sobre $\underline C_\Lambda$. Las condiciones
siguientes son equivalentes:

(i) $\varphi$ es liso;

(ii) todo morfismo sobreyectivo
$\beta:b'\longrightarrow\varphi(a)$ en $\underline B$ puede
completarse en un diagrama $(*)$ en el que
$\varphi(a')\longrightarrow b'$ es un
$\underline C_\Lambda$-isomorfismo;

(iii) toda pequeña prolongación
$\beta:b'\longrightarrow\varphi(a)$ en $\underline B$ puede
completarse en un diagrama conmutativo $(*)$.

\underline{Observación 2.2(b).} Para todo grupoide cofibrado
$\underline A$ sobre $\underline C_\Lambda$, el funtor natural
\[
 \underline A\longrightarrow[\underline A]
\]
es mínimamente liso.

\underline{Observación 2.2(c).} Para cada
$R\in\widehat{\underline C}_\Lambda$, sea
\[
 h_R:\underline C_\Lambda\longrightarrow(\mathrm{Ens})
\]
el funtor definido por
\[
 h_R(S)=\operatorname{Hom}_{\Lambda\text{-alg}}(R,S)
 \qquad(S\in\underline C_\Lambda).
\]
Este determina un grupoide discreto cofibrado $\underline h_R$ sobre
$\underline C_\Lambda$. Un morfismo $R\longrightarrow S$ en
$\widehat{\underline C}_\Lambda$ induce un funtor
\[
 \underline h_S\longrightarrow\underline h_R.
\]
Este funtor es liso si y solo si $S$ es un anillo de series formales
sobre $R$.

%%%%%%%%%%  scan idx 62  |  folio 51  %%%%%%%%%%

\underline{Observación 2.2(d).} Dado
$R\in\widehat{\underline C}_\Lambda$, denotemos por
$(R,\underline C_\Lambda)$ la categoría de las flechas de
$\widehat{\underline C}_\Lambda$ con fuente $R$ y blanco en
$\underline C_\Lambda$. Sea $\underline A$ un grupoide cofibrado
sobre $\underline C_\Lambda$. Dado
$\xi\in\operatorname{Ob}\widehat{\underline A}$ con $P(\xi)=R$,
sea
\[
 \underline A'\longrightarrow(R,\underline C_\Lambda)
\]
el producto fibrado de
$\underline A\longrightarrow\underline C_\Lambda$ por el funtor
canónico
$(R,\underline C_\Lambda)\longrightarrow\underline C_\Lambda$.
Elijamos una sección $\xi^\#$ tal que
$\xi^\#(1_R)=\xi$. Así, para todo morfismo
$f:R\longrightarrow S$, elegimos un $f$-morfismo
$\xi\longrightarrow\xi(f)$.

Obtenemos un funtor
\[
 \xi^\#:\underline h_R\longrightarrow\underline A
\]
sobre $\underline C_\Lambda$ enviando
$(R\xrightarrow{f}S)$ a $\xi(f)$. Salvo equivalencia natural, este
funtor no depende de la elección de $\xi^\#$.

La afirmación de que
$\xi^\#:\underline h_R\longrightarrow\underline A$ es
(mínimamente) liso significa que todo morfismo sobreyectivo
$a'\xrightarrow{\alpha}\xi(f)$ en $\underline A$ puede completarse
en un diagrama conmutativo
\[
\begin{tikzcd}[row sep=large,column sep=huge]
 \xi(f') \arrow[r] \arrow[dr,"{\xi(\,)}"']
 & a' \arrow[d,"\alpha"]\\
 & \xi(f)
\end{tikzcd}
\]
(y, en el caso mínimo, que
$\underline h_R(k[\varepsilon])\longrightarrow
\underline A(k[\varepsilon])$ es biyectivo). Esto significa
precisamente que $\xi$ es un objeto formal (mínimamente) versal de
$\underline A$ sobre $\underline C_\Lambda$.

Así, la existencia de un objeto formal (mínimamente) versal de
$\underline A$ equivale a la existencia de un funtor (mínimamente)
liso
\[
 \underline h_R\longrightarrow\underline A
\]
sobre $\underline C_\Lambda$ para algún
$R\in\operatorname{Ob}(\widehat{\underline C}_\Lambda)$.

